\documentclass[12pt]{article}
\usepackage{amsmath, amssymb, geometry, graphicx, hyperref, color}
\geometry{letterpaper, margin=1in}

\title{Homework 16.3 Fundamental Theorem of Line Integrals}
\author{Tashfeen Omran}
\date{April 2026}

\begin{document}

\maketitle
\tableofcontents
\newpage

\section{Part 1: Comprehensive Introduction, Context, and Prerequisites}

\subsection{Core Concepts}
The topic of this chapter is the \textbf{Fundamental Theorem for Line Integrals (FTLI)}, a profound extension of the Fundamental Theorem of Calculus into multivariable spaces. 

When dealing with a vector field $\mathbf{F}(x, y) = P(x, y)\mathbf{i} + Q(x, y)\mathbf{j}$, a line integral $\int_C \mathbf{F} \cdot d\mathbf{r}$ calculates the total ``work'' done by the field along a curve $C$. Generally, this work depends on the specific path taken from point $A$ to point $B$. However, for a special class of vector fields called \textbf{conservative vector fields}, the work done depends \textit{only} on the starting and ending points, rendering the integral \textbf{path-independent}. 

A vector field $\mathbf{F}$ is conservative if it is the gradient of some scalar function $f$, meaning $\mathbf{F} = \nabla f$. The function $f$ is known as the \textbf{scalar potential function}. If such a function exists, evaluating the line integral becomes a simple matter of evaluating $f$ at the curve's endpoints, bypassing the need for tedious parameterization entirely.

\subsection{Intuition and Derivation}
Recall the Fundamental Theorem of Calculus (Part II) from single-variable calculus: 
\[
\int_{a}^{b} f'(x) \, dx = f(b) - f(a)
\]
This tells us that the integral of a rate of change over an interval is simply the total accumulated change from the start to the end.

The FTLI is the exact same principle in higher dimensions. If $\mathbf{F} = \nabla f$, then $\mathbf{F}$ represents the multi-dimensional rate of change (gradient) of the potential function $f$. By the Chain Rule for paths $\mathbf{r}(t)$, the dot product $\nabla f(\mathbf{r}(t)) \cdot \mathbf{r}'(t)$ is exactly the derivative of $f$ with respect to time $t$. Integrating this rate of change along the path exactly cancels out the differential, leaving only the net change in the potential $f$ from the start to the end.

\subsection{Historical Context and Motivation}
The concept of conservative fields and potential functions originated in 19th-century physics, specifically in classical mechanics and electromagnetism. Physicists like George Green, Carl Friedrich Gauss, and Lord Kelvin were trying to model gravitational and electrostatic forces. They discovered that moving a mass in a gravitational field, or a charge in an electrostatic field, in a closed loop required zero net work. This implied the forces were derived from a scalar ``potential'' (e.g., gravitational potential energy or voltage). The mathematical machinery developed to formalize this physics intuition gave birth to the Fundamental Theorem of Line Integrals.

\subsection{Key Formulas}
\begin{enumerate}
    \item \textbf{Fundamental Theorem for Line Integrals:}
    If $C$ is a piecewise-smooth curve given by $\mathbf{r}(t)$ for $a \le t \le b$, and $f$ is a differentiable function whose gradient vector $\nabla f$ is continuous on $C$, then:
    \[
    \int_C \nabla f \cdot d\mathbf{r} = f(\mathbf{r}(b)) - f(\mathbf{r}(a))
    \]
    
    \item \textbf{Component Test for Conservative Fields (2D):}
    Let $\mathbf{F}(x, y) = P(x, y)\mathbf{i} + Q(x, y)\mathbf{j}$. If the partial derivatives are continuous on an open, simply-connected region, then $\mathbf{F}$ is conservative if and only if:
    \[
    \frac{\partial P}{\partial y} = \frac{\partial Q}{\partial x}
    \]

    \item \textbf{Component Test for Conservative Fields (3D):}
    For $\mathbf{F}(x, y, z) = P\mathbf{i} + Q\mathbf{j} + R\mathbf{k}$, the field is conservative if its curl is zero ($\nabla \times \mathbf{F} = \mathbf{0}$), which translates to:
    \[
    \frac{\partial R}{\partial y} = \frac{\partial Q}{\partial z}, \quad \frac{\partial P}{\partial z} = \frac{\partial R}{\partial x}, \quad \frac{\partial Q}{\partial x} = \frac{\partial P}{\partial y}
    \]
\end{enumerate}

\subsection{Prerequisites}
To master this topic, you must be proficient in:
\begin{itemize}
    \item \textbf{Partial Differentiation:} Taking derivatives with respect to one variable while treating others as constants.
    \item \textbf{Partial Integration:} Integrating a multivariable expression with respect to one variable. Crucially, the ``constant of integration'' becomes a \textit{function} of the remaining variables (e.g., integrating with respect to $x$ yields $+ g(y)$).
    \item \textbf{Basic Vectors:} Understanding dot products and vector magnitude.
\end{itemize}

\newpage
\section{Part 2: Detailed Homework Solutions}

\subsection{Problem 1: Example Verifications}
\textbf{Determine whether or not the given vector field is conservative.}

\subsubsection*{Part (a)}
\textbf{Given:} $\mathbf{F}(x, y) = (5x - 6y)\mathbf{i} + (2x - 3)\mathbf{j}$

\textbf{Solution Setup:}
Let $P(x, y) = 5x - 6y$ and $Q(x, y) = 2x - 3$. We must check if $\frac{\partial P}{\partial y} = \frac{\partial Q}{\partial x}$.

\textbf{Step 1: Compute partial derivatives.}
\[
\frac{\partial P}{\partial y} = \frac{\partial}{\partial y}(5x - 6y) = -6
\]
\[
\frac{\partial Q}{\partial x} = \frac{\partial}{\partial x}(2x - 3) = 2
\]

\textbf{Step 2: Compare and conclude.}
Since $-6 \neq 2$, we have $\frac{\partial P}{\partial y} \neq \frac{\partial Q}{\partial x}$. 

\textbf{Final Answers for the boxes:}
\begin{itemize}
    \item $\frac{\partial P}{\partial y} = \mathbf{-6}$
    \item $\frac{\partial Q}{\partial x} = \mathbf{2}$
    \item Since $\frac{\partial P}{\partial y} \neq \frac{\partial Q}{\partial x}$, $\mathbf{F}$ is \textbf{not conservative}.
\end{itemize}

\subsubsection*{Part (b)}
\textbf{Given:} $\mathbf{F}(x, y) = (3 + 4xy)\mathbf{i} + (2x^2 - 7y^2)\mathbf{j}$

\textbf{Solution Setup:}
Let $P(x, y) = 3 + 4xy$ and $Q(x, y) = 2x^2 - 7y^2$.

\textbf{Step 1: Compute partial derivatives.}
\[
\frac{\partial P}{\partial y} = \frac{\partial}{\partial y}(3 + 4xy) = 4x
\]
\[
\frac{\partial Q}{\partial x} = \frac{\partial}{\partial x}(2x^2 - 7y^2) = 4x
\]

\textbf{Step 2: Compare and conclude.}
Since $4x = 4x$, the partials are equal. The domain is $\mathbb{R}^2$ which is open and simply connected.

\textbf{Final Answers for the boxes:}
\begin{itemize}
    \item $\frac{\partial P}{\partial y} = \mathbf{4x}$
    \item $\frac{\partial Q}{\partial x} = \mathbf{4x}$
    \item Therefore $\mathbf{F}$ is \textbf{conservative}.
\end{itemize}

---

\subsection{Problem 2: Conservative Check and Potential}
\textbf{Determine whether or not $\mathbf{F}$ is a conservative vector field. If it is, find a function $f$ such that $\mathbf{F} = \nabla f$. (If the vector field is not conservative, enter DNE.)}

\textbf{Given:} $\mathbf{F}(x, y) = (xy + y^5)\mathbf{i} + (x^2 + 5xy)\mathbf{j}$

\textbf{Step 1: Identify components and test for conservatism.}
\[
P(x, y) = xy + y^5 \quad \implies \quad \frac{\partial P}{\partial y} = x + 5y^4
\]
\[
Q(x, y) = x^2 + 5xy \quad \implies \quad \frac{\partial Q}{\partial x} = 2x + 5y
\]

\textbf{Step 2: Compare.}
Since $x + 5y^4 \neq 2x + 5y$, the field is not conservative. Therefore, no such potential function exists.

\textbf{Final Answer:} $f(x,y) = $ \textbf{DNE}

---

\subsection{Problem 3: Finding a Scalar Potential}
\textbf{Given:} $\mathbf{F}(x, y) = (y^2 - 4x)\mathbf{i} + 2xy\mathbf{j}$

\textbf{Step 1: Test for conservatism.}
\[
P(x,y) = y^2 - 4x \implies \frac{\partial P}{\partial y} = 2y
\]
\[
Q(x,y) = 2xy \implies \frac{\partial Q}{\partial x} = 2y
\]
The field is conservative.

\textbf{Step 2: Set up the potential function equations.}
We know $f_x = y^2 - 4x$ and $f_y = 2xy$.

\textbf{Step 3: Integrate $f_x$ with respect to $x$.}
\[
f(x, y) = \int (y^2 - 4x) \, dx = xy^2 - 2x^2 + g(y)
\]

\textbf{Step 4: Differentiate with respect to $y$ and equate to $Q$.}
\[
\frac{\partial}{\partial y}[xy^2 - 2x^2 + g(y)] = 2xy + g'(y)
\]
We are given $f_y = 2xy$, so:
\[
2xy + g'(y) = 2xy \quad \implies \quad g'(y) = 0 \quad \implies \quad g(y) = K
\]

\textbf{Final Answer:} $f(x, y) = \mathbf{xy^2 - 2x^2}$

---

\subsection{Problem 4: Transcendentals Potential}
\textbf{Given:} $\mathbf{F}(x, y) = (ye^x + \sin(y))\mathbf{i} + (e^x + x\cos(y))\mathbf{j}$

\textbf{Step 1: Test for conservatism.}
\[
P(x,y) = ye^x + \sin(y) \implies \frac{\partial P}{\partial y} = e^x + \cos(y)
\]
\[
Q(x,y) = e^x + x\cos(y) \implies \frac{\partial Q}{\partial x} = e^x + \cos(y)
\]
The field is conservative.

\textbf{Step 2: Integrate $P$ with respect to $x$.}
\[
f(x, y) = \int (ye^x + \sin(y)) \, dx = ye^x + x\sin(y) + g(y)
\]

\textbf{Step 3: Differentiate with respect to $y$ and solve for $g(y)$.}
\[
f_y = \frac{\partial}{\partial y}[ye^x + x\sin(y) + g(y)] = e^x + x\cos(y) + g'(y)
\]
Equating this to $Q(x, y)$:
\[
e^x + x\cos(y) + g'(y) = e^x + x\cos(y) \quad \implies \quad g'(y) = 0 \quad \implies \quad g(y) = K
\]

\textbf{Final Answer:} $f(x, y) = \mathbf{ye^x + x\sin(y)}$

---

\subsection{Problem 5: Trigonometric Potential}
\textbf{Given:} $\mathbf{F}(x, y) = (y^8\cos(x) + \cos(y))\mathbf{i} + (8y^7\sin(x) - x\sin(y))\mathbf{j}$

\textbf{Step 1: Test for conservatism.}
\[
P_y = 8y^7\cos(x) - \sin(y)
\]
\[
Q_x = 8y^7\cos(x) - \sin(y)
\]
The field is conservative.

\textbf{Step 2: Integrate $P$ with respect to $x$.}
\[
f(x, y) = \int (y^8\cos(x) + \cos(y)) \, dx = y^8\sin(x) + x\cos(y) + g(y)
\]

\textbf{Step 3: Differentiate with respect to $y$ and match with $Q$.}
\[
f_y = 8y^7\sin(x) - x\sin(y) + g'(y)
\]
Equating to $Q$:
\[
8y^7\sin(x) - x\sin(y) + g'(y) = 8y^7\sin(x) - x\sin(y) \quad \implies \quad g'(y) = 0
\]

\textbf{Final Answer:} $f(x, y) = \mathbf{y^8\sin(x) + x\cos(y)}$

---

\subsection{Problem 6: Rational Functions Potential}
\textbf{Given:} $\mathbf{F}(x, y) = \left(\ln(y) + \frac{y}{x}\right)\mathbf{i} + \left(\ln(x) + \frac{x}{y}\right)\mathbf{j}$

\textbf{Step 1: Test for conservatism.}
\[
P_y = \frac{\partial}{\partial y}\left(\ln(y) + yx^{-1}\right) = \frac{1}{y} + \frac{1}{x}
\]
\[
Q_x = \frac{\partial}{\partial x}\left(\ln(x) + xy^{-1}\right) = \frac{1}{x} + \frac{1}{y}
\]
The field is conservative (assuming $x, y > 0$).

\textbf{Step 2: Integrate $P$ with respect to $x$.}
\[
f(x, y) = \int \left(\ln(y) + \frac{y}{x}\right) \, dx = x\ln(y) + y\ln(x) + g(y)
\]

\textbf{Step 3: Differentiate with respect to $y$.}
\[
f_y = \frac{x}{y} + \ln(x) + g'(y)
\]
Equate to $Q$:
\[
\frac{x}{y} + \ln(x) + g'(y) = \ln(x) + \frac{x}{y} \quad \implies \quad g'(y) = 0
\]

\textbf{Final Answer:} $f(x, y) = \mathbf{x\ln(y) + y\ln(x)}$

---

\subsection{Problem 7: Evaluating Line Integral using FTLI}
\textbf{Consider $\mathbf{F}$ and $C$ below. $\mathbf{F}(x, y) = \langle 2x, 8y \rangle$, $C$ is the arc of the parabola $x = y^2$ from $(4, -2)$ to $(1, 1)$.}

\subsubsection*{Part (a): Find a function $f$ such that $\mathbf{F} = \nabla f$}
\textbf{Step 1: Integrate components.}
\[
P = 2x \implies f = \int 2x \, dx = x^2 + g(y)
\]
\textbf{Step 2: Find $g(y)$.}
\[
f_y = g'(y) = Q = 8y \implies g(y) = \int 8y \, dy = 4y^2
\]
\textbf{Final Answer for (a):} $f(x,y) = \mathbf{x^2 + 4y^2}$

\subsubsection*{Part (b): Use part (a) to evaluate $\int_C \mathbf{F} \cdot d\mathbf{r}$}
By the Fundamental Theorem of Line Integrals:
\[
\int_C \mathbf{F} \cdot d\mathbf{r} = f(B) - f(A) = f(1, 1) - f(4, -2)
\]
\textbf{Step 1: Evaluate $f$ at the endpoints.}
\[
f(1, 1) = (1)^2 + 4(1)^2 = 1 + 4 = 5
\]
\[
f(4, -2) = (4)^2 + 4(-2)^2 = 16 + 4(4) = 16 + 16 = 32
\]
\textbf{Step 2: Subtract.}
\[
5 - 32 = -27
\]
\textbf{Final Answer for (b):} $\mathbf{-27}$

---

\subsection{Problem 8: Evaluating Line Integral of Rational Function}
\textbf{Given:} $\mathbf{F}(x, y) = (5 + 8xy^2)\mathbf{i} + 8x^2y\mathbf{j}$, $C$ is arc of hyperbola $y = 1/x$ from $(1, 1)$ to $(5, 1/5)$.

\subsubsection*{Part (a): Find a function $f$ such that $\mathbf{F} = \nabla f$}
\textbf{Step 1: Integrate $P$ with respect to $x$.}
\[
f(x, y) = \int (5 + 8xy^2) \, dx = 5x + 4x^2y^2 + g(y)
\]
\textbf{Step 2: Differentiate with respect to $y$ and solve for $g(y)$.}
\[
f_y = 8x^2y + g'(y)
\]
Since $Q = 8x^2y$, we have $g'(y) = 0$.

\textbf{Final Answer for (a):} $f(x,y) = \mathbf{5x + 4x^2y^2}$

\subsubsection*{Part (b): Use part (a) to evaluate $\int_C \mathbf{F} \cdot d\mathbf{r}$}
\[
\int_C \mathbf{F} \cdot d\mathbf{r} = f(5, 1/5) - f(1, 1)
\]
\textbf{Step 1: Evaluate potential at endpoints.}
\[
f\left(5, \frac{1}{5}\right) = 5(5) + 4(5)^2\left(\frac{1}{5}\right)^2 = 25 + 4(25)\left(\frac{1}{25}\right) = 25 + 4 = 29
\]
\[
f(1, 1) = 5(1) + 4(1)^2(1)^2 = 5 + 4 = 9
\]
\textbf{Step 2: Subtract.}
\[
29 - 9 = 20
\]
\textbf{Final Answer for (b):} $\mathbf{20}$

---

\subsection{Problem 9: Finding a 3D Potential Function}
\textbf{Given:} $\mathbf{F}(x, y, z) = 4y^2\mathbf{i} + (8xy + 3e^{4z})\mathbf{j} + 12ye^{4z}\mathbf{k}$

We are guided through a step-by-step framework to find $f$.

\textbf{Step 1: State the partial derivatives from the vector field.}
\[
\text{(1)} \quad f_x(x, y, z) = \mathbf{4y^2}
\]
\[
\text{(2)} \quad f_y(x, y, z) = \mathbf{8xy + 3e^{4z}}
\]
\[
\text{(3)} \quad f_z(x, y, z) = \mathbf{12ye^{4z}}
\]

\textbf{Step 2: Integrating (1) with respect to $x$.}
\[
\text{(4)} \quad f(x, y, z) = \int 4y^2 \, dx = \mathbf{4xy^2} + g(y, z)
\]

\textbf{Step 3: Differentiate (4) wrt $y$ and compare with (2).}
\[
f_y(x, y, z) = \mathbf{8xy} + g_y(y, z)
\]
Comparing this to (2), which is $8xy + 3e^{4z}$, yields:
\[
g_y(y, z) = \mathbf{3e^{4z}}
\]

\textbf{Step 4: Integrate $g_y$ wrt $y$.}
Thus, $g(y, z) = 3ye^{4z} + h(z)$ and we write (4) as:
\[
f(x, y, z) = \mathbf{4xy^2} + 3ye^{4z} + h(z)
\]

\textbf{Step 5: Differentiate wrt $z$ and compare with (3).}
\[
f_z = 12ye^{4z} + h'(z)
\]
Comparing with (3), $12ye^{4z}$, we obtain:
\[
h'(z) = \mathbf{0}
\]
Therefore $h(z) = K$.

\textbf{Step 6: Give the desired function.}
\[
f(x, y, z) = \mathbf{4xy^2 + 3ye^{4z}} + K
\]
\textit{Note: The final box expects $4xy^2 + 3ye^{4z}$.}

\newpage
\section{Part 3: In-Depth Analysis of Problems and Techniques}

\subsection{A) Problem Types and General Approach}
The homework assignment features three main categories of problems:
\begin{enumerate}
    \item \textbf{Verification of Conservative Fields (Problems 1, 2):}
    \textit{Approach:} Extract $P$ and $Q$. Differentiate $P$ with respect to $y$, and $Q$ with respect to $x$. Check if they are strictly equal. If they are, the field is conservative. If they differ even by a negative sign, it is not.
    \item \textbf{Finding 2D and 3D Potential Functions (Problems 3-6, 9):}
    \textit{Approach:} Use a systematic partial integration technique. Integrate $P$ with respect to $x$, appending an unknown function of the remaining variables (e.g., $g(y)$). Differentiate this temporary function with respect to the next variable, equate it to the next component of the vector field, and integrate again to discover the unknown function. 
    \item \textbf{Evaluating Line Integrals using FTLI (Problems 7, 8):}
    \textit{Approach:} Instead of parameterizing the complex curve $C$ (like the arc of a hyperbola), first verify the field is conservative and find its potential function $f$. Then, simply plug the terminal point and starting point into $f$ and subtract. 
\end{enumerate}

\subsection{B) Key Algebraic and Calculus Manipulations}
\begin{itemize}
    \item \textbf{Partial Integration Constants:} 
    Unlike single-variable calculus where integration yields a numerical constant $+ C$, integrating an expression like $\int 4y^2 \, dx$ in a multivariable context yields $4xy^2 + g(y, z)$. The ``constant'' must account for all variables treated as fixed during the integration step. This was perfectly demonstrated in \textbf{Problem 9}.
    \item \textbf{Derivative/Integral Matching Strategy:}
    In \textbf{Problems 4 and 5}, mixing transcendental functions (like $e^x\sin(y)$) can look intimidating to integrate. However, the systematic matching strategy sidesteps complex integration. By integrating the simpler component first and comparing the partial derivative of the result against the remaining component, the unknown function $g(y)$ often collapses to $0$.
\end{itemize}

\newpage
\section{Part 4: ``Cheatsheet'' and Tips for Success}

\subsection{Summary of Most Important Formulas}
\begin{itemize}
    \item \textbf{2D Conservative Check:} $\frac{\partial P}{\partial y} = \frac{\partial Q}{\partial x}$
    \item \textbf{3D Conservative Check (Curl):} $\nabla \times \mathbf{F} = \mathbf{0}$
    \item \textbf{FTLI:} $\int_C \mathbf{F} \cdot d\mathbf{r} = f(\text{end point}) - f(\text{start point})$
\end{itemize}

\subsection{Tricks and Shortcuts}
\begin{itemize}
    \item \textbf{The Symmetry Shortcut for Potential Functions:} When looking for $f(x,y)$, you can integrate $P$ with respect to $x$, and separately integrate $Q$ with respect to $y$. Look at the distinct terms generated by both integrals. The total potential function $f$ is simply the sum of all \textit{unique} terms from both integrals! 
    \textit{Example (Problem 3):} $P = y^2 - 4x \implies \int P \, dx = xy^2 - 2x^2$. $Q = 2xy \implies \int Q \, dy = xy^2$. The unique terms are $xy^2$ and $-2x^2$. Thus, $f = xy^2 - 2x^2$.
    \item \textbf{Skip Parameterization immediately:} If a problem asks you to evaluate a line integral but gives you a very messy curve (e.g., $y = 1/x$ or $x = y^2$), that is an immediate red flag that you should be using the FTLI. Do not attempt to compute $\mathbf{r}'(t)$ unless finding a potential function fails.
\end{itemize}

\subsection{Common Pitfalls}
\begin{itemize}
    \item \textbf{Mixing up the Cross-Partials:} A classic mistake is checking $\frac{\partial P}{\partial x} = \frac{\partial Q}{\partial y}$. This is entirely incorrect. Always cross the variables: differentiate the $x$-component ($P$) with respect to $y$, and the $y$-component ($Q$) with respect to $x$.
    \item \textbf{Dropping the function constant $g(y)$:} If you forget $g(y)$ and just write $+ C$, you will fail to capture remaining terms that depend purely on $y$.
    \item \textbf{Assuming Conservative blindly:} You cannot use $f(B) - f(A)$ if the field is not conservative. Always verify first.
\end{itemize}

\newpage
\section{Part 5: Conceptual Synthesis and The ``Big Picture''}

\subsection{A) Thematic Connections}
The central theme here is \textbf{Path Independence}. In our previous studies, the path taken drastically changed the resulting value (e.g., Arc Length is highly dependent on the wiggles in the path). Here, we encounter special environments (conservative fields) where the universe only ``cares'' about your net displacement, not the journey. 

This mirrors the concept of \textbf{State Variables} in physics and thermodynamics. An object's temperature or altitude is a state variable; the change in altitude only depends on the start and end heights, regardless of the stairs climbed to get there.

\subsection{B) Forward and Backward Links}
\textbf{Backward Link:} This is the culmination of the single-variable FTC (Part 2) and partial derivatives. We are lifting the 1D rule $\int f' = \Delta f$ into any arbitrary curve looping through 3D space.

\textbf{Forward Link:} This lays the critical foundation for \textbf{Green's Theorem} and \textbf{Stokes' Theorem}. Soon, we will deal with fields that are \textit{not} conservative. Green's and Stokes' theorems measure exactly \textit{how far} a field is from being conservative (using Curl) and relate the line integral around a closed loop to a double integral over the enclosed surface.

\newpage
\section{Part 6: Real-World Application and Modeling (Finance Focus)}

\subsection{A) Concrete Scenarios}
\begin{enumerate}
    \item \textbf{Arbitrage-Free Asset Pricing:} In quantitative finance and stochastic calculus, the price of a derivative (like a European option) depends on state variables, such as the current stock price $S$ and time $t$. A portfolio constructed to hedge against these changes will have sensitivities (the "Greeks" like Delta $\Delta$ and Theta $\Theta$). For the market to be free of arbitrage, trading a sequence of assets in a closed loop must result in a net cost/payoff of strictly zero. This implies the "vector field" of state-prices must be a conservative vector field, and the derivative's price is the scalar potential function!
    \item \textbf{Path-Dependent vs. Path-Independent Exotic Options:} An Asian Option's payout depends on the average price of the underlying asset over a period of time (the specific path taken matters). However, a standard European Call Option depends only on the asset's price at expiration (path-independent). Evaluating the expected payoff of a European option under a risk-neutral measure can be mathematically framed as evaluating a potential function at a terminal state, mirroring the FTLI.
\end{enumerate}

\subsection{B) Model Problem Setup}
\textbf{Scenario:} Arbitrage Pricing Model
Let $x$ represent the price of Stock A, and $y$ represent the price of Stock B. The cost to restructure a financial portfolio continuously in response to price changes is given by the vector field of required cash flows:
\[
\mathbf{F}(x, y) = P(x, y)\mathbf{i} + Q(x, y)\mathbf{j}
\]
where $P$ is the local cash sensitivity to Stock A, and $Q$ is the local cash sensitivity to Stock B.

An automated trading algorithm buys and sells these stocks, tracing a curve $C$ in the price-space $(x, y)$ from Monday's opening prices $(x_0, y_0)$ to Friday's closing prices $(x_1, y_1)$. The total cost of trading over the week is $\int_C \mathbf{F} \cdot d\mathbf{r}$. 

If the market is strictly arbitrage-free, then the total cost to transition from $(x_0, y_0)$ to $(x_1, y_1)$ cannot depend on the daily price fluctuations (the path $C$), but only on the start and end prices. Therefore, the financial engineers must verify the no-arbitrage condition:
\[
\frac{\partial P}{\partial y} = \frac{\partial Q}{\partial x}
\]
If true, they find the absolute pricing function $V(x, y)$ (the potential function) such that $\nabla V = \mathbf{F}$, and the total cost is simply $V(x_1, y_1) - V(x_0, y_0)$.

\newpage
\section{Part 7: Common Variations and Untested Concepts}

The homework provided excellent coverage of 2D potentials and the standard FTLI. However, a couple of standard curriculum concepts were missing.

\subsection{Missing Concept 1: The Closed Loop Property}
If a field is conservative, the line integral over any \textit{closed loop} (where the start point and end point are the same) is exactly zero.
\[
\oint_C \mathbf{F} \cdot d\mathbf{r} = f(A) - f(A) = 0
\]
\textbf{Example Problem:} Evaluate $\oint_C (2x)\mathbf{i} + (2y)\mathbf{j} \cdot d\mathbf{r}$ where $C$ is the unit circle $x^2 + y^2 = 1$.
\textbf{Solution:} Notice $P_y = 0$ and $Q_x = 0$. The field is conservative. Since the unit circle is a closed loop starting and ending at $(1,0)$, the FTLI guarantees the integral is exactly $0$. No parameterization necessary!

\subsection{Missing Concept 2: 3D Curl Check for Conservatism}
Before finding a 3D potential function (as in Homework Problem 9), you usually must prove the field is conservative by showing the Curl is the zero vector: $\nabla \times \mathbf{F} = \mathbf{0}$.
\textbf{Example Problem:} Determine if $\mathbf{F} = \langle yz, xz, xy \rangle$ is conservative.
\textbf{Solution:} We check the three cross-partials:
1) $R_y = Q_z \implies \frac{\partial}{\partial y}(xy) = x$, $\frac{\partial}{\partial z}(xz) = x$. ($x = x$)
2) $P_z = R_x \implies \frac{\partial}{\partial z}(yz) = y$, $\frac{\partial}{\partial x}(xy) = y$. ($y = y$)
3) $Q_x = P_y \implies \frac{\partial}{\partial x}(xz) = z$, $\frac{\partial}{\partial y}(yz) = z$. ($z = z$)
Since all match, the curl is zero, and $\mathbf{F}$ is conservative.

\newpage
\section{Part 8: Advanced Diagnostic Testing: ``Find the Flaw''}

Below are five solutions containing subtle, common errors. Your task as the grader is to spot the mistake, explain why it is wrong in one sentence, and provide the correct solution.

\subsection{Flawed Problem 1}
\textbf{Problem:} Evaluate the line integral of $\mathbf{F}(x, y) = \langle y^2, x \rangle$ along the curve $C$ from $(0,0)$ to $(1,1)$.
\textbf{Flawed Solution:} 
1. We use the FTLI. We need $\nabla f = \mathbf{F}$.
2. $\int y^2 \, dx = xy^2$.
3. Thus $f(x, y) = xy^2$.
4. $\int_C \mathbf{F} \cdot d\mathbf{r} = f(1,1) - f(0,0) = 1(1)^2 - 0 = 1$.

\subsection{Flawed Problem 2}
\textbf{Problem:} Find a potential function for $\mathbf{F}(x, y) = \langle e^y, xe^y \rangle$.
\textbf{Flawed Solution:}
1. $f_x = e^y \implies f = \int e^y \, dx = xe^y + C$.
2. We check $f_y$: $\frac{\partial}{\partial y}(xe^y + C) = xe^y$, which matches $Q$.
3. Final Answer: $f(x, y) = xe^y + C$.

\subsection{Flawed Problem 3}
\textbf{Problem:} Determine if $\mathbf{F}(x, y) = \langle x^2y, xy^2 \rangle$ is conservative.
\textbf{Flawed Solution:}
1. Let $P = x^2y$ and $Q = xy^2$.
2. We test the partials: $\frac{\partial P}{\partial x} = 2xy$ and $\frac{\partial Q}{\partial y} = 2xy$.
3. Since $2xy = 2xy$, the field is conservative.

\subsection{Flawed Problem 4}
\textbf{Problem:} Evaluate $\int_C \mathbf{F} \cdot d\mathbf{r}$ for $\mathbf{F}(x, y) = \langle 2x, 2y \rangle$, where $C$ is the upper half of the circle $x^2 + y^2 = 4$ traversed from $(2,0)$ to $(-2,0)$.
\textbf{Flawed Solution:}
1. The field is conservative with potential $f(x, y) = x^2 + y^2$.
2. The curve is a circle, which means it returns to the $x$-axis. By the property of conservative fields, the integral over loops is $0$.
3. Final answer is $0$.

\subsection{Flawed Problem 5}
\textbf{Problem:} Find the potential function for the 3D field $\mathbf{F}(x, y, z) = \langle z, z, x + y \rangle$.
\textbf{Flawed Solution:}
1. $f_x = z \implies f = xz + g(z)$.
2. $f_y = z \implies f = yz + h(z)$.
3. Adding them up gives $f(x, y, z) = xz + yz$. Checking with $R$: $f_z = x + y$, which matches.
4. Final Answer: $f(x, y, z) = xz + yz$.

---
\vspace{1cm}
\textbf{Answer Key for Grader (You):}
\begin{itemize}
    \item \textbf{Flaw 1:} The student failed to test if the field is conservative; $P_y = 2y$ and $Q_x = 1$, so it is not conservative and FTLI cannot be used. (Must parameterize path to solve).
    \item \textbf{Flaw 2:} The student integrated with respect to $x$ and added a numerical constant $+ C$ instead of an arbitrary function of $y$, $+ g(y)$, which represents a critical conceptual misunderstanding of partial integration.
    \item \textbf{Flaw 3:} The student took the partial derivatives with respect to the \textit{wrong variables} ($P_x$ and $Q_y$ instead of $P_y$ and $Q_x$); $P_y = x^2$ and $Q_x = y^2$, so it is not conservative.
    \item \textbf{Flaw 4:} The student assumed a half-circle is a closed loop, but the endpoints $(2,0)$ and $(-2,0)$ are different. (Correct evaluation is $f(-2,0) - f(2,0) = 4 - 4 = 0$, but the reasoning of a closed loop is completely flawed).
    \item \textbf{Flaw 5:} The student incorrectly added the partial potential functions together rather than recognizing $g(y,z)$ from the first step and correctly setting up $f = xz + g(y,z)$. 
\end{itemize}

\end{document}